\documentclass[11pt,a4paper]{article}
\usepackage[T1]{fontenc}
\usepackage[utf8]{inputenc}
\usepackage{lmodern,amsmath,amssymb}
\usepackage[margin=25mm]{geometry}
\usepackage{longtable,booktabs,array,calc}
\usepackage[hidelinks]{hyperref}
\hypersetup{pdftitle={Continuous-time expansion for Kogut--Susskind: $SU(3)$ gapped for $g^2\ge388$, gap approac},pdfauthor={navstokgap research working notes}}
\newcommand{\tightlist}{\setlength{\itemsep}{0pt}\setlength{\parskip}{0pt}}
\setlength{\emergencystretch}{3em}
\setcounter{secnumdepth}{0}
\begin{document}
\hypertarget{continuous-time-expansion-for-kogutsusskind-su3-gapped-for-g2ge388-gap-approaching-83g2hbar-ca}{%
\section{\texorpdfstring{Continuous-time expansion for Kogut--Susskind:
\(SU(3)\) gapped for \(g^2\ge388\), gap approaching
\((8/3)g^2\,\hbar c/a\)}{Continuous-time expansion for Kogut--Susskind: SU(3) gapped for g\^{}2\textbackslash ge388, gap approaching (8/3)g\^{}2\textbackslash,\textbackslash hbar c/a}}\label{continuous-time-expansion-for-kogutsusskind-su3-gapped-for-g2ge388-gap-approaching-83g2hbar-ca}}

Yarotsky's theorem gives the Kogut--Susskind gap for \(SU(3)\) only
beyond \(g^2\sim10^{101}\), because its time discretization pays
\(e^{64t_0}\) per perturbed site. The same statement with usable numbers
comes from the Duhamel expansion of \(e^{-tH}\) in the plaquette term
with the electric evolution kept exact between insertions: a polymer gas
in continuous time whose weights carry the exact decay
\(e^{-E\,\Delta\tau}\) of every excited link, so no neighbourhood is
ever left unpenalized. Three structural facts do the work. The electric
term is diagonal in the character basis with eigenvalue
\(\varepsilon\sum_\ell C_2(r_\ell)/C_2(f)\),
\(\varepsilon=(g^2/2)C_2(f)=2g^2/3\). The plaquette operator
\(w_p=\operatorname{tr}U_p+\operatorname{tr}U_p^\dagger\) acts on four
links with \(\|w_p\|\le2N=6\), and
\(V=(2/g^2)\sum_p(N-\operatorname{Re}\operatorname{tr}U_p)\) is
\(-(1/g^2)\sum_pw_p\) up to a constant. And Gauss's law forces every
intermediate excited set to be a closed flux loop, so at least four
links are excited whenever any are, \(E\ge4\varepsilon\) on the
orthogonal complement of the vacuum. The Kotecký--Preiss criterion, in
its measure-space form, with the weight
\(a(\gamma)=n(\gamma)+\lambda|\operatorname{supp}\gamma|\) and
\(\lambda=\varepsilon(1-\theta)\), then gives a gap of at least
\(4\lambda=(8/3)g^2(1-\theta)\) in units \(\hbar c/a\) whenever
\[g^4\ \ge\ \max\Big(\frac{3132}{\theta},\ \frac{783}{\sqrt{\theta(1-\theta)}}\Big),\]
for the class of polymers in which each insertion touches the current
excited set; so \(\Delta_{\rm KS}\ge\frac43g^2\,\hbar c/a\) for
\(g^2\ge79\), \(\Delta_{\rm KS}\ge2.4\,g^2\,\hbar c/a\) for
\(g^2\ge177\), and the bound approaches the flux-loop energy as
\(g\to\infty\), which fixes the constant \(\gamma\) of the T2 theorem at
\(4\). The general polymers, in which a later insertion may start a
separate piece that is bridged afterwards, are counted explicitly by a
spanning tree of consecutive touchings, at a factor \(24\) per
insertion, which makes the rigorous threshold \(g^2\ge388\) for the same
\(\frac43g^2\) gap. Constants explicit; nothing promoted.

\hypertarget{the-expansion}{%
\subsection{1. The expansion}\label{the-expansion}}

Units \(\hbar c/a=1\). Write \(H=H_E-W+{\rm const}\) with
\(H_E=(g^2/2)\sum_\ell(-\Delta_\ell)\) and \(W=(1/g^2)\sum_pw_p\). In
the Peter--Weyl basis \(H_E\) is diagonal, its eigenvalue on a link
configuration \(\sigma=(r_\ell)\) being
\(E(\sigma)=(g^2/2)\sum_\ell C_2(r_\ell)\ge\varepsilon\,|S(\sigma)|\),
with \(S(\sigma)\) the set of links carrying a nontrivial representation
and \(\varepsilon=(g^2/2)C_2(f)\), since \(C_2(f)=4/3\) is the least
nonzero Casimir of \(SU(3)\). Duhamel's formula gives
\[e^{-tH}=\sum_{n\ge0}\int_{0<\tau_1<\dots<\tau_n<t}
e^{-(t-\tau_n)H_E}\,W\,e^{-(\tau_n-\tau_{n-1})H_E}\cdots W\,e^{-\tau_1H_E},\]
all terms positive. Insert, between consecutive factors, the resolution
\(1=\sum_SP_S\) with \(P_S\) the projection onto configurations whose
excited set is exactly \(S\); each \(P_S\) is a product over links. A
\textbf{history} is then a sequence \((p_k,\tau_k,S_k)_{k=1}^n\) with
\(S_0=\emptyset\), and its weight is
\[\Big\langle\Omega_0,\prod_{k}\big(P_{S_k}\tfrac1{g^2}w_{p_k}P_{S_{k-1}}\big)e^{-\Delta\tau_kH_E}\ \Omega_0\Big\rangle,
\qquad\Big|\cdot\Big|\ \le\ \Big(\frac{2N}{g^2}\Big)^n\prod_{k=1}^{n-1}e^{-\varepsilon|S_k|\Delta\tau_k},\]
because \(\|w_p\|\le2N\) and \(e^{-\Delta\tau H_E}P_S\) has norm at most
\(e^{-\varepsilon|S|\Delta\tau}\).

Two constraints on the \(S_k\). First,
\(\operatorname{links}(p_k)\subseteq S_{k-1}\cup S_k\): acting on a
trivial link, \(w_p\) makes it nontrivial (\(f\otimes1=f\)), so a link
of \(p_k\) can be trivial afterwards only if it was nontrivial before.
Second, \textbf{Gauss's law}: \(\Omega_0\) and every \(w_p\) are gauge
invariant, so every intermediate configuration is, and the
representations on the links at a vertex fuse to the trivial one; a
nonempty excited set is therefore a union of closed flux loops and
\(|S_k|\ge4\) whenever \(S_k\ne\emptyset\). This is where the
non-abelian theory pays nothing extra: the same constraint holds for
\(U(1)\).

\textbf{Factorization.} The support of a history is
\(\operatorname{supp}\gamma=\bigcup_k S_k\times[\tau_k,\tau_{k+1}]\subset\mathcal E\times[0,t]\).
If a history splits into pieces with disjoint supports, its weight is
the product of the pieces' weights, because \(P_S\),
\(e^{-\Delta\tau H_E}\) and \(\Omega_0\) factor over links and each
\(w_p\) acts on four links; and the ordered time integral over the whole
history is the product of the ordered integrals over the pieces. So
\[Z(t)=\langle\Omega_0,e^{-tH}\Omega_0\rangle=\sum_{\{\gamma_i\}\ \text{compatible}}\prod_iw(\gamma_i),\]
a gas of \textbf{polymers}, the connected histories, two polymers being
compatible when their supports are disjoint. Between consecutive
insertions of a connected polymer the excited set is nonempty, or the
support would split in time; hence \(|S_k|\ge4\) for \(1\le k\le n-1\).

\hypertarget{the-convergence-criterion}{%
\subsection{2. The convergence
criterion}\label{the-convergence-criterion}}

The Kotecký--Preiss criterion (Commun. Math. Phys. 103 (1986) 491) in
its form for polymers indexed by a measure space (Poghosyan and
Ueltschi, J. Math. Phys. 50 (2009) 053509; metadata level) is: if for
every \(\gamma\)
\[\int_{\gamma'\not\sim\gamma}|w(\gamma')|\,e^{a(\gamma')}\,d\gamma'\ \le\ a(\gamma),\]
then the cluster expansion of \(\log Z\) converges absolutely and
uniformly in the volume, and truncated correlations obey the standard
bounds. Here \(d\gamma'\) is the sum over shapes \((p_k,S_k)_{k=1}^{n}\)
times \(d\tau_1\prod_{k=1}^{n-1}d\Delta\tau_k\).

Take \[a(\gamma)=n(\gamma)+\lambda\,|\operatorname{supp}\gamma|,\qquad
|\operatorname{supp}\gamma|=\sum_{k=1}^{n-1}|S_k|\Delta\tau_k,\qquad 0<\lambda<\varepsilon .\]
Then
\(e^{a(\gamma')}|w(\gamma')|\le(2Ne/g^2)^{n'}\prod_ke^{-(\varepsilon-\lambda)|S_k'|\Delta\tau_k'}\).

\emph{Overlap.} \(\gamma'\not\sim\gamma\) means some link \(\ell\) is
excited in both at a common time. For a fixed shape and internal times
of \(\gamma'\), the set of \(\tau_1'\) producing an overlap through link
\(\ell\in S_k'\) during the interval in which \(\gamma\) excites
\(\ell\) has measure at most \(T_\ell(\gamma)+\Delta\tau_k'\), where
\(T_\ell(\gamma)\) is the total time \(\gamma\) keeps \(\ell\) excited,
and \(\sum_\ell T_\ell(\gamma)=|\operatorname{supp}\gamma|\).

\emph{Integrals.}
\(\int_0^\infty e^{-(\varepsilon-\lambda)|S|\Delta\tau}d\Delta\tau=\frac1{(\varepsilon-\lambda)|S|}\)
and
\(\int_0^\infty\Delta\tau\,e^{-(\varepsilon-\lambda)|S|\Delta\tau}d\Delta\tau=\frac1{((\varepsilon-\lambda)|S|)^2}\le\frac1{4(\varepsilon-\lambda)}\cdot\frac1{(\varepsilon-\lambda)|S|}\).

Collecting, with
\[Q_\lambda\ :=\ \sum_{\text{shapes}\ni\ell_0}(n-1)\Big(\frac{2Ne}{g^2}\Big)^n\prod_{k=1}^{n-1}\frac{1}{(\varepsilon-\lambda)|S_k|}\]
the weighted number of polymer shapes through a fixed link, the left
side of the criterion is at most
\(Q_\lambda|\operatorname{supp}\gamma|+\frac{Q_\lambda}{4(\varepsilon-\lambda)}\,n(\gamma)\),
and the criterion holds whenever
\[Q_\lambda\ \le\ \lambda\qquad\text{and}\qquad Q_\lambda\le4(\varepsilon-\lambda).\]

\hypertarget{counting-shapes}{%
\subsection{3. Counting shapes}\label{counting-shapes}}

\textbf{Adjacent growth.} Consider first the shapes in which every
insertion \(p_k\), \(k\ge2\), shares a link with the current excited set
\(S_{k-1}\). The first plaquette contains \(\ell_0\): in three
dimensions a link lies in four plaquettes, so \(4\) choices, and
\(S_1=\operatorname{links}(p_1)\) is forced. At step \(k\ge2\) there are
at most \(4|S_{k-1}|\) plaquettes touching \(S_{k-1}\) and at most
\(2^4=16\) choices of \(S_k\). Each step therefore contributes at most
\[\frac{2Ne}{g^2}\cdot\frac{1}{(\varepsilon-\lambda)|S_{k-1}|}\cdot64|S_{k-1}|
=\frac{128Ne}{g^2(\varepsilon-\lambda)}\ =:\ u,\] the size \(|S_{k-1}|\)
cancelling between the entropy and the energy denominator, which is the
mechanism that keeps the count geometric. So
\[Q_\lambda^{\rm adj}\ \le\ 4\,\frac{2Ne}{g^2}\sum_{m\ge1}m\,u^m=\frac{8Ne}{g^2}\,\frac{u}{(1-u)^2}.\]

\textbf{General shapes, counted by a tree.} A later insertion may start
a piece disjoint from the current support that is bridged afterwards. To
count all connected shapes, build the graph \(G'\) on the insertions in
which, for every link \(\ell\) and every maximal interval during which
\(\ell\) is excited, consecutive insertions touching \(\ell\) are joined
by an edge. Connectedness of the support is connectedness of \(G'\), and
the edges attached to a given link tile its excitation intervals without
overlap, so for any spanning tree \(T\subseteq G'\)
\[\sum_{e\in T}|\Delta\tau_e|\ \le\ |\operatorname{supp}\gamma|,\qquad\text{hence}\qquad
e^{-(\varepsilon-\lambda)|\operatorname{supp}\gamma|}\le\prod_{e\in T}e^{-(\varepsilon-\lambda)|\Delta\tau_e|}.\]
Each insertion has at most eight slots for \(G'\)-edges, the previous
and next toucher of each of its four links, so \(T\) is a rooted tree
with children in labelled slots: at most \(4^n\) plane shapes times
\(8^{n-1}\) slot assignments. Given the parent and the slot, the child
plaquette contains the slot's link and differs from the parent, at most
\(3\) choices; the excited set after it, at most \(16\); the insertion
factor \(2Ne/g^2\); and the edge time integral
\(\int_0^\infty e^{-(\varepsilon-\lambda)\Delta\tau}d\Delta\tau=1/(\varepsilon-\lambda)\),
the time direction being fixed by the slot. The per-vertex factor is
\[u_{\rm gen}=32\cdot3\cdot16\cdot\frac{2Ne}{g^2}\cdot\frac1{\varepsilon-\lambda}
=\frac{3072Ne}{g^2(\varepsilon-\lambda)}=24\,u ,\] and the sum over all
connected shapes through a fixed link is bounded by the adjacent-growth
sum with \(u\) replaced by \(24u\); the overlap factor of Section 2
loses its \(1/4\) because only one link's decay is used per edge, so the
second condition reads \(Q\le\varepsilon-\lambda\). The loss against
adjacent growth is the factor \(24\), coming from the tree entropy
\(32\), the \(3\) plaquettes per slot and the single-link decay in place
of \(1/(4(\varepsilon-\lambda))\). Both counts are given below; the
general one is rigorous as written, the adjacent one is what a careful
count of the same shapes should approach.

\hypertarget{numbers-for-su3}{%
\subsection{\texorpdfstring{4. Numbers for
\(SU(3)\)}{4. Numbers for SU(3)}}\label{numbers-for-su3}}

With \(N=3\), \(\varepsilon=2g^2/3\) and
\(\lambda=\varepsilon(1-\theta)\),
\[u=\frac{128\cdot3e}{g^2\cdot\theta\cdot2g^2/3}=\frac{576e}{\theta g^4}=\frac{1566}{\theta g^4}.\]
Impose \(u\le\frac12\) and, for the first condition of Section 2,
\(\frac{8Ne}{g^2}\cdot4u\le\lambda\), that is
\(96e\cdot\frac{1566}{\theta g^4}\le\frac23g^2(1-\theta)\cdot g^2/g^2\);
cleaning up,
\[g^4\ \ge\ \frac{3132}{\theta}\qquad\text{and}\qquad g^8\ \ge\ \frac{6.13\times10^5}{\theta(1-\theta)},\]
the second condition of Section 2 being weaker. Hence:

\begin{longtable}[]{@{}lll@{}}
\toprule\noalign{}
\(\theta\) & threshold \(g^2\) & gap
\(\ge4\lambda=\frac83g^2(1-\theta)\) \\
\midrule\noalign{}
\endhead
\bottomrule\noalign{}
\endlastfoot
\(1/2\) & \(79\) & \(\frac43g^2\) \\
\(1/4\) & \(112\) & \(2g^2\) \\
\(1/10\) & \(177\) & \(2.4\,g^2\) \\
\(\theta\to0\) & \(g^2\gtrsim\sqrt{3132/\theta}\) & \(\to\frac83g^2\) \\
\end{longtable}

With the general count, \(u_{\rm gen}=37584/(\theta g^4)\), the
conditions \(u_{\rm gen}\le\frac12\) and \(4Fu_{\rm gen}\le\lambda\)
give \(g^4\ge75168/\theta\) and
\(g^8\ge1.47\times10^7/(\theta(1-\theta))\):

\begin{longtable}[]{@{}lll@{}}
\toprule\noalign{}
\(\theta\) & rigorous threshold \(g^2\) & gap
\(\ge\frac83g^2(1-\theta)\) \\
\midrule\noalign{}
\endhead
\bottomrule\noalign{}
\endlastfoot
\(1/2\) & \(388\) & \(\frac43g^2\) \\
\(1/4\) & \(548\) & \(2g^2\) \\
\(1/10\) & \(867\) & \(2.4\,g^2\) \\
\end{longtable}

The gap here is read off as in
\href{wilson-strong-coupling-explicit.md}{the Wilson note} §3: running
the criterion with \(a(\gamma)=n+\lambda|\operatorname{supp}\gamma|\)
puts a factor \(e^{-\lambda|\operatorname{supp}X|}\) on every cluster, a
cluster spanning a time interval of length \(t\) has
\(|\operatorname{supp}|\ge4t\) by Gauss's law, so vacuum-subtracted
correlations in Euclidean time decay at rate at least \(4\lambda\), and
for a positive semigroup that is a lower bound on the gap on the cyclic
subspace of local observables, which in finite volume is everything by
T1.

\textbf{Consequence for T2.} The T2 theorem's constant \(\gamma\),
defined by \(\Delta\ge\gamma(g^2/2)C_2\), is \(4(1-\theta)\) in this
expansion and tends to \(4\): the gap is the energy of one flux loop,
four links in the fundamental, up to corrections of relative order
\(1/g^4\). This agrees with the two-sided strong-coupling statement of
\href{su3-constants.md}{the SU(3) constants} §2, whose upper side is
\(12g^2\), and closes the ratio \(18/\gamma\) there to \(4.5\).

\hypertarget{why-this-beats-the-time-discretized-argument}{%
\subsection{5. Why this beats the time-discretized
argument}\label{why-this-beats-the-time-discretized-argument}}

Yarotsky's Lemma 3 gives up the decay of every link in the range of a
perturbed site for a whole step \(t_0\), and \(t_0\ge7\) is forced by
the entropy of the space-time lattice; the product is
\(e^{64t_0}\sim e^{450}\) per insertion. In continuous time each excited
link decays exactly for exactly as long as it is excited, the insertion
carries no step, and the only entropy is the choice of the next
plaquette and excited set, which the energy denominator
\(1/((\varepsilon-\lambda)|S_k|)\) pays for link by link. The threshold
drops from \(10^{101}\) to \(10^2\) for the same model and the same
conclusion.

\hypertarget{consequence-for-state}{%
\subsection{6. Consequence for STATE}\label{consequence-for-state}}

The T2 theorem for the Kogut--Susskind Hamiltonian now has an explicit
threshold, \(g^2\ge388\) for \(SU(3)\) rigorously and \(g^2\ge79\) in
the adjacent-growth class, and an explicit gap approaching the flux-loop
energy \((8/3)g^2\,\hbar c/a\). Together with the Wilson transfer-matrix
result, the strong-coupling boundary of the intermediate region is of
order \(g^2\sim10^2\) for both standard regularizations. The weak side
remains the crude \(1/g^2\gtrsim10^2\), and the intermediate region
between them is unchanged.
\end{document}
